\ifdefined\ICLRLayout\clearpage\fi
\subsection{Theoretical Analysis}
\label{sec:theory}

The analysis establishes sufficient conditions under which context aggregation
improves Base and characterizes how task-gated fusion controls prediction
changes according to context support. Let
$g_0$ and $g_{\mathrm{final}}$ return $o_i^{(0)}$ and
$o_i^{\mathrm{final}}$, respectively, and define
$\mathcal R_t(g)=\mathbb E_{(x,y)\sim P_t}[\ell(g(x),y)]$. For a loss $\ell$
that is convex in its prediction argument, define the improvements of context
$r$ and the context aggregate over Base as
\begin{equation}
  \Delta_i^{(r)}
  =\ell(o_i^{(0)},y_i^t)-\ell(o_i^{(r)},y_i^t),
  \qquad
  \Delta_i
  =\ell(o_i^{(0)},y_i^t)-\ell(\bar o_i,y_i^t).
  \label{eq:context_improvements}
\end{equation}
These quantities use target labels only for analysis, never for context
construction, support scoring, or prediction.

\paragraph{Complementary context coverage.}
For a change family $h$, alternative contexts improve Base on different
subsets of queries. We use $C_h(\Phi)$ to denote the target probability that
at least one context in bank $\Phi$ improves Base by at least $\delta_+$. A new
context is complementary when it adds beneficial queries not already covered
by the bank. Formal definitions and the
diminishing-returns property are given in Appendix~\ref{app:repair-coverage}.

\begin{theorem}[Risk bound for weighted context aggregation]
\label{thm:effective_context_repair}
For query $i$ and fixed bank $\Phi$, suppose that a nonempty effective set
$\mathcal G_i$ satisfies, for
some $\delta_+>0$, $\delta_-\geq0$, and $\rho_+>\rho_-$,
\begin{equation}
\begin{aligned}
\Delta_i^{(r)}&\geq\delta_+,&\rho_i^{(r)}(\Phi)&\geq\rho_+&& (r\in\mathcal G_i),\\
\Delta_i^{(r)}&\geq-\delta_-,&\rho_i^{(r)}(\Phi)&\leq\rho_-&& (r\in\mathcal I(\Phi)\setminus\mathcal G_i).
\end{aligned}
\label{eq:repair_conditions}
\end{equation}
Let $\Pi_i=\sum_{r\in\mathcal G_i}\pi_r$ and
$W_i=\sum_{r\in\mathcal G_i}w_i^{(r)}$.  Then
\begin{equation}
\begin{aligned}
W_i&\geq\underline W_i
=\left[1+\frac{1-\Pi_i}{\Pi_i}
e^{-(\rho_+-\rho_-)}\right]^{-1},\\
\Delta_i&\geq\gamma_i
:=\underline W_i\delta_+-(1-\underline W_i)\delta_-.
\end{aligned}
\label{eq:repair_margin}
\end{equation}
Hence $\underline W_i>\delta_-/(\delta_++\delta_-)$ implies
$\gamma_i>0$ and
\begin{equation}
  \ell(o_i^{\mathrm{final}},y_i^t)-\ell(o_i^{(0)},y_i^t)
  \leq-\alpha_i\Delta_i
  \leq-\alpha_{\max}U_iE_i\gamma_i.
  \label{eq:query_repair_gain}
\end{equation}

Let $\mathcal A_\Phi$ contain the queries for which the preceding separation
conditions hold and $\gamma_i>0$.  The target-risk difference obeys
\begin{equation}
  \mathcal R_t(g_{\mathrm{final}})-\mathcal R_t(g_0)
  \leq-\Gamma_t(\Phi)+\varepsilon_t(\Phi),
  \label{eq:changed_risk_gain}
\end{equation}
where
\begin{equation}
\begin{aligned}
\Gamma_t(\Phi)
&:=\alpha_{\max}\mathbb E_{P_t}
  [U_iE_i\gamma_i\mathbf 1\{i\in\mathcal A_\Phi\}],\\
\varepsilon_t(\Phi)
&:=\alpha_{\max}\mathbb E_{P_t}
  [U_iE_i(-\Delta_i)_+\mathbf 1\{i\notin\mathcal A_\Phi\}].
\end{aligned}
\label{eq:repair_residual_terms}
\end{equation}
Therefore $\Gamma_t(\Phi)>\varepsilon_t(\Phi)$ implies
$\mathcal R_t(g_{\mathrm{final}})<\mathcal R_t(g_0)$.
\end{theorem}

\paragraph{Proof sketch.}
Score separation yields $W_i\geq\underline W_i$ for the normalized weight on
$\mathcal G_i$. Convexity then converts this bound into the aggregate
improvement $\gamma_i$, which yields the query-level gain in
Eq.~\eqref{eq:query_repair_gain}. Splitting the target expectation over
$\mathcal A_\Phi$ and its complement gives
Eq.~\eqref{eq:changed_risk_gain}. The complete proof appears in
Appendix~\ref{app:theorem-proof}.

\paragraph{Implications.}
A larger effective prior mass $\Pi_i$ or a wider score gap
$\rho_+-\rho_-$ raises $\underline W_i$ and hence the guaranteed aggregate
gain $\gamma_i$. Complementary contexts determine potential beneficial coverage,
while score weighting and gating determine how that coverage affects realized
risk. The following exact identity gives the magnitude of each task-specific
update relative to Base:
\begin{equation}
  \lVert o_i^{\mathrm{final}}-o_i^{(0)}\rVert
  =\alpha_{\max}U_iE_i
   \lVert\bar o_i-o_i^{(0)}\rVert.
  \label{eq:base_preservation_identity}
\end{equation}
Thus, the prediction recovers Base whenever $U_iE_i=0$. Task-specific
classification and regression bounds are given in Appendix~\ref{app:theory}.
